% Propietario conservado: /Users/ruben/Documents/New project/output/EXTREMA_MEDIA_RAZON_RECIPROCIDAD_ES_EN_20260918/ARTICULO/ES/source/sections/realizacion_lc.tex
\subsection{Réalisation inductive--capacitive de la forme d’action}
\label{x_reciprocidad:nuclear:realizacion-lc}

La forme d’action déjà construite admet une réalisation électrique dans
une carte munie d’une horloge \(\omega>0\) et d’une impédance de
référence \(Z_*>0\). Ces deux grandeurs conservent leurs unités et
accompagnent le paramètre de forme \(\eta\). On définit
\begin{equation}
 L_\eta=\frac{Z_*}{\omega}e^{-\eta},\qquad
 C_\eta=\frac{e^\eta}{Z_*\omega},\qquad
 Z_\eta=\sqrt{L_\eta/C_\eta}.
 \label{x_reciprocidad:hmtII:eq:ii-elipse-lc}
\end{equation}
Ici, \(C_\eta\) désigne la capacité électrique, distincte du contre-angle \(C^*\).

\begin{proposition}[Conservation du produit dynamique]
\label{x_reciprocidad:nuclear:lc-producto}
La réalisation \eqref{x_reciprocidad:hmtII:eq:ii-elipse-lc} satisfait
\begin{equation}
 L_\eta C_\eta=\omega^{-2},\qquad
 \frac{Z_\eta}{Z_*}=e^{-\eta}=\frac{b_S}{a_S},
 \label{x_reciprocidad:hmtII:eq:ii-elipse-impedancia}
\end{equation}
où \(a_S=\sqrt{2\hbar}\,e^{\eta/2}\) et
\(b_S=\sqrt{2\hbar}\,e^{-\eta/2}\) sont les demi-axes de
\cref{x_reciprocidad:iv:ellipse:semiejes}.
Avec \(Q=\sqrt{Z_*}q\) et \(P=\Phi/\sqrt{Z_*}\), l’énergie prend la forme
\[
 H_{LC}=\frac{q^2}{2C_\eta}+\frac{\Phi^2}{2L_\eta}
       =\frac\omega2(e^{-\eta}Q^2+e^\eta P^2).
\]
Sur \((Q,P)=(a_S\cos\theta,b_S\sin\theta)\), les contributions
électrique et magnétique sont \(\hbar\omega\cos^2\theta\) et
\(\hbar\omega\sin^2\theta\) ; l’énergie totale est \(\hbar\omega\).
\end{proposition}
\begin{proof}
La multiplication de \(L_\eta\) et \(C_\eta\) annule les facteurs
\(Z_*\) et \(e^{\pm\eta}\), tandis que leur quotient est
\(Z_*^2e^{-2\eta}\). La racine positive prouve l’identité d’impédance.
Substituer \(q=Q/\sqrt{Z_*}\) et \(\Phi=\sqrt{Z_*}P\) donne la forme
de \(H_{LC}\). Les carrés des demi-axes annulent les facteurs
de forme restants ; la somme utilise \(\cos^2\theta+\sin^2\theta=1\).
\end{proof}

La convention hamiltonienne
\(\dot Q=\partial_PH\), \(\dot P=-\partial_QH\) attribue
\(\dot\theta=-\omega\) à ce paramétrage. La forme, l’horloge
et la carte dimensionnelle restent différenciées. La réponse
constitutive du vide et l’impédance \(Z_\eta\) sont reliées
par l’inversion des canaux de \cref{x_reciprocidad:iii:prop:forma-LC} ; chacune
conserve sa définition et sa normalisation.

